\section{Two clearance marks separated by actual reflections}
\label{sec:xf-jacobi}
The same-flight theorem of Section~\ref{sec:dc-two-mark} is an
inherited result. We now treat marks on \emph{different} flights.
The intervening reflections, rather than an independence assumption,
supply the new transversality. This section proves finite-itinerary
geometric identities on the original physical billiard. The uniform
sum in the next section explicitly uses the retained physical pressure
input; neither section asserts a normalized roof estimate.

\subsection{Physical feet and the reflection derivative}
Put $\kappa_-=1/R_+$ and $\kappa_+=1/R_-$. All radii, clearances,
collision coordinates and measures are those of the preceding sections.
The following strengthening of the inherited interior-foot bound is
useful because it is independent of every visited incidence.
\begin{lemma}[A uniform physical foot standoff]
\label{lem:xf-standoff}
For a genuine mark of clearance $0<u\le\delta_*=3/200$ on a
physical flight, its perpendicular foot has distance at least
\begin{equation}\label{eq:xf-standoff}
 a_\circ=\sqrt{(1-R_+)^2-(R_++\delta_*)^2}
          =\frac{\sqrt{1827}}{200}>\tau_-=\frac{3}{50}
\end{equation}
from each incident collision. In particular the foot is strictly
interior, uniformly for $R\in[9/20,47/100]$.
\end{lemma}
\begin{proof}
The interior assertion is already Lemma~\ref{lem:cf-foot}; it uses
that a third center cannot realize its small clearance at an incident
endpoint. If $q$ is either incident contact and $z$ is the marked
center, lattice spacing gives $|z-q|\ge1-R$. Orthogonal decomposition
along the actual flight gives $t^2=|z-q|^2-(R+u)^2$ for the distance
$t>0$ between that endpoint and the foot. The right side is at least
$(1-R_+)^2-(R_++\delta_*)^2=1827/40000$. Finally
$1827>144$ proves $a_\circ>3/50$. No extra collision at the foot
has been introduced.
\end{proof}

Consider a first mark on flight $0$ and a second on flight $h$, with
$h\ge1$. Let $q_1,\ldots,q_h$ be the actual intervening collision
points, $c_i>0$ their normal incidences, $a$ the distance from the first
foot to $q_1$, and $b$ that from $q_h$ to the second foot. Write
$\ell_i=|q_{i+1}-q_i|$ for $1\le i<h$. Thus
$a,b\ge a_\circ$ and $\ell_i\ge\tau_-$. Define
\begin{align}
 P(t)&=\begin{pmatrix}1&t\\0&1\end{pmatrix},&
 K_i&=\begin{pmatrix}1&0\\2/(Rc_i)&1\end{pmatrix},\notag\\
 \mathcal M_h
   &=P(b)K_hP(\ell_{h-1})K_{h-1}\cdots
                           P(\ell_1)K_1P(a)
     =\begin{pmatrix}A_h&B_h\\C_h&D_h\end{pmatrix}.
                                      \label{eq:xf-product}
\end{align}
For $h=1$ the middle product is empty:
$\mathcal M_1=P(b)K_1P(a)$. The letter $B_h$ here is a matrix
entry; it is not the outer localization parameter $B$.

\begin{theorem}[The exact cross-flight clearance Jacobian]
\label{thm:xf-jacobian}
Translate the first marked center to $0$ and parametrize its line by
$(\vartheta,u_1,\sigma_1)$ as in Theorem~\ref{thm:cf-flux}.
On every regular chart where the intervening physical itinerary and
second marked center are fixed, let $b_2$ be the second signed impact
parameter. At fixed $u_1$,
\begin{equation}\label{eq:xf-jacobian}
 \left|\partial_\vartheta b_2\right|=B_h>0,
 \qquad
 \left|\det D_{(\vartheta,u_1)}(u_1,u_2)\right|=B_h,
 \quad u_2=|b_2|-R.
\end{equation}
Moreover $\det\mathcal M_h=1$, and the pointwise lower bounds
\begin{align}
 B_h&\ge a+b+\sum_{i=1}^{h-1}\ell_i,\label{eq:xf-length-lower}\\
 B_h&\ge a_\circ(1+2\kappa_-\tau_-)^{h-1}
                                      \label{eq:xf-exponential-lower}
\end{align}
hold for every such trajectory. They require no lower cutoff on any
$c_i$. A visited incidence tending to zero makes its positive kick
larger, not smaller.
\end{theorem}
\begin{proof}
For a differentiable variation of an oriented line let $J$ denote its
transverse displacement and $J'$ its direction variation, measured at
the reference point on that line. Sliding the reference point along the
same line adds only the usual free-flight term. Hence travel by distance
$t$ sends $(J,J')$ to $(J+tJ',J')$.

Here is the reflection calculation with its signs. With outward normal
$n$, tangent $\mathcal J n$ and curvature $\kappa=1/R$, write
$v^-=-cn+p\mathcal Jn$ and $v^+=cn+p\mathcal Jn$, where
$\mathcal J$ is rotation by $\pi/2$. A contact variation is
$\delta q=\mathcal Jn\,\delta s$. Its transverse components in the
incoming and outgoing direction frames are $-c\delta s$ and
$c\delta s$, respectively. Specular reflection gives
$\delta\vartheta^+=2\kappa\delta s-\delta\vartheta^-$.
Reverse the transverse frame after each reflection, applying the same
sign reversal to $J'$ as to $J$. The contact displacement then becomes
continuous, while
\[
              J'^+=J'^-+\frac{2\kappa}{c}J.
\]
Thus the exact linearized reflection in these unfolded transverse
frames is $K_i$. This is a derivative of the genuine reflection rule,
not an auxiliary reflecting wall at a marked disk.

The first foot is $-\sigma_1(R+u_1)v_\vartheta^\perp$. Differentiating
it at fixed $u_1$ gives a vector parallel to $v_\vartheta$; its
transverse component is zero and its direction variation is one.
Consequently the transverse field begins with $(J,J')=(0,1)$.
Propagation to the second foot gives $(B_h,D_h)$. Differentiating the
second signed distance gives the negative of this transverse
component, up to the $h$ frame reversals. For example, at the last
contact, differentiation of
$b_2=(z_2-q_h)\cdot v_h^\perp$ gives
$-\delta q_h\cdot v_h^\perp-b\,\delta\vartheta_h$.
Its magnitude is exactly $B_h$. Multiplication by the second impact
sign has unit magnitude, proving both identities in
\eqref{eq:xf-jacobian}.

Each factor in \eqref{eq:xf-product} has determinant one and
nonnegative entries. Replacing all $K_i$ by the identity can only
decrease the product entrywise; the resulting upper-right entry is
the sum of the free lengths. This proves
\eqref{eq:xf-length-lower}. At the first collision $J=a\ge a_\circ$.
At each subsequent collision, propagation from its predecessor obeys
$J_{i+1}=J_i+\ell_iJ_i'^+$ with
$J_i'^+\ge 2\kappa_-J_i$. Hence
$J_{i+1}\ge(1+2\kappa_-\tau_-)J_i$. The final free flight also
increases $J$, which proves \eqref{eq:xf-exponential-lower}.
All assertions have been proved on regular physical charts; a
one-sided limiting matrix with an infinite grazing kick is not
called a finite derivative of a singular physical return.
\end{proof}

\subsection{Global multiplicity, including disconnected word domains}
A local nonzero Jacobian does not bound the number of its roots.
The following global argument supplies precisely that missing fact.
A \emph{tangentially augmented itinerary} consists of the first marked
center, the ordered $h$ actual intervening disk centers, the second
marked center, and the two signed impacts. The marked centers may
coincide with one another or with centers visited at other times.
They are not the incident centers of their own marked flights.

\begin{theorem}[One trajectory per tangentially augmented itinerary]
\label{thm:xf-unique}
Fix $h\ge1$, the augmented itinerary, $R$, and $u_1,u_2$ in their
original clearance ranges. There is at most one regular physical
broken path between the two feet realizing these data, including its
oriented line at each endpoint. This is a global statement: no
connectedness of its admissible word domain is assumed.
\end{theorem}
\begin{proof}
For two candidate configurations
$Q=(q_0,q_1,\ldots,q_h,q_{h+1})$ and $\widehat Q$, the endpoint
positions lie on the two auxiliary circles of radii $R+u_1$ and
$R+u_2$, and the other positions lie on the actual radius-$R$ disks.
These auxiliary circles only locate the feet; no new collision is
counted. The Cartesian broken length
$\mathcal S(Q)=\sum_{i=0}^{h}|q_{i+1}-q_i|$ is convex and
continuously differentiable at each candidate because every segment
has positive length. Gradient monotonicity gives
\begin{equation}\label{eq:xf-convex}
 (\nabla\mathcal S(Q)-\nabla\mathcal S(\widehat Q))
                         \cdot(Q-\widehat Q)\ge0.
\end{equation}
At an actual interior collision
$\nabla_i\mathcal S=-2c_i n_i$, so its contribution on the left is
$-R(c_i+\widehat c_i)|n_i-\widehat n_i|^2$.
At a marked endpoint the velocity is tangent to its auxiliary circle,
with its oriented tangent fixed by the prescribed impact sign.
Thus the endpoint gradient is a fixed sign times $\mathcal Jn_i$.
Its difference is orthogonal to the contact difference, and its
contribution to \eqref{eq:xf-convex} is zero. Summing shows that
every actual interior contact agrees, since $c_i,\widehat c_i>0$.

For completeness, agreement of the contacts alone is not used to
skip the endpoint argument. Write the two convex support gaps for
$\mathcal S$ and add them. Their sum is the left side of
\eqref{eq:xf-convex}, now zero. Each gap is a sum of nonnegative
support gaps of the Euclidean norm for individual flight vectors.
Equality in the norm support inequality, with nonzero reference
vector, forces the two corresponding vectors to be positive
multiples of one another. Therefore all corresponding flight
directions agree. At either endpoint its direction and prescribed
impact sign determine its auxiliary-circle normal uniquely, hence
the foot itself. This proves global uniqueness. Repeated disk centers
cause no identification of contact variables at different times and
do not affect the argument.
\end{proof}

\subsection{An exact physical two-clearance change of variables}
Set
\begin{equation}\label{eq:xf-candidates}
 \mathcal Z_\circ=\{z\in\Lambda\setminus\{0\}:
           |z|\le\tau_++2R_++\delta_*\},\qquad
 N_\circ=|\mathcal Z_\circ|.
\end{equation}
This is the same finite end-candidate geometry as the one-mark
construction. After translating the first marked center to zero,
the first actual collision center has at most $N_\circ$ choices;
the second marked center relative to the last actual collision
center has at most $N_\circ$ choices. The intervening physical word
has length $h-1$, \emph{not} $h$. It is empty when $h=1$.
Let $\mathcal A_{h,R}(u_1,u_2)$ be the set of these augmented
itineraries with a regular physical realization, each counted once.
Denote the associated original collision state and the two marked
centers in its lifted initial frame by $(x_a,z_{1,a},z_{2,a})$.

\begin{theorem}[One-copy coarea across arbitrarily many true reflections]
\label{thm:xf-coarea}
For any $h\ge1$ and every nonnegative Borel weight $f$ of the initial
physical state and the two marks,
\begin{align}
 &\int \sum_{z_1\in\mathcal Z_{\delta_1}(x)}
          \sum_{z_2\in\mathcal Z_{\delta_2}(T_R^hx)}
                  f(x,z_1,z_2)\,d\nu(x)\notag\\
 &\quad=\frac1{4\pi R}\int_0^{\delta_1}\int_0^{\delta_2}
       \sum_{a\in\mathcal A_{h,R}(u_1,u_2)}
                 \frac{f(x_a,z_{1,a},z_{2,a})}{B_h(x_a)}
                                    \,du_2\,du_1.\label{eq:xf-coarea}
\end{align}
In the left side, the center on flight $h$ is translated into the
same original lifted frame as the first. The two marks are distinct
as flight--center pairs even when their lattice centers agree.
All original return, displacement, collision, physical-roof and
section tests may be included in $f$ without altering the identity.
\end{theorem}
\begin{proof}
Apply the exact first-mark formula of Theorem~\ref{thm:cf-flux};
its prefactor is $(4\pi R)^{-1}$ and its parameters are
$(\vartheta,u_1,\sigma_1)$. On each regular physical branch the
second center has a unique impact sign. The second substitution has
Jacobian $B_h$ by Theorem~\ref{thm:xf-jacobian}. The two signs at
each end are included among the augmented itineraries; no extra
factor of two, radius, incidence, collision count, or word count is
permitted. Theorem~\ref{thm:xf-unique} ensures that overlapping
charts of a single augmented itinerary cannot produce different roots.

To justify the global substitution, exhaust the regular parameter
space by countably many relatively compact inverse-function charts,
and disjointize their overlaps before summation. On each chart use
the ordinary change-of-variables theorem first for bounded
nonnegative weights. The chart exhaustion and then monotone
approximation give \eqref{eq:xf-coarea} for arbitrary $f\ge0$.
Finite-horizon geometry gives finitely many physical words at each
fixed $h$, although this number is never used as a bound. The usual
singular collision set is null for collision area and hence for
first-mark area on every chart. Its image under a regular local
change of variables is null as well. Values on those bulk-null
sets therefore do not alter this identity. This argument makes no
claim that their one-sided paired endpoint traces vanish. Membership
in $\mathcal A_{h,R}$ retains all physical flights and feet; enlarging
that set is allowed only for subsequent positive upper bounds.
\end{proof}
